\section{Clustering Methodology}
\label{sec:clustering}

We evaluate four user-pairing methods for PD\mbox{-}NOMA: \emph{Static}, \emph{Balanced}, \emph{Blossom} (general maximum-weight matching), and \emph{Bipartite proportional-fair (PF)} matching.  
All methods pre-enforce the SIC guard before admitting a pair; any unpaired users are served via OMA fallback.

\subsection{Static Strong–Weak Pairing}
Users are sorted in ascending order of $|h|$, then paired symmetrically:
\[
\mathcal{P}_{\text{static}}=\{(u_{(i)},u_{(N+1-i)})\mid i=1,\ldots,N/2\}.
\]
This ensures each weak user is coupled with a strong one, maximizing disparity for SIC feasibility.  
As shown in Fig.~\ref{fig:static}, the static clustering pattern connects users at opposite ends of the channel-gain spectrum, producing long cross-cell links that emphasize strong–weak diversity.  
While this approach maximizes channel disparity, it may lead to reduced fairness among mid-gain users.

\begin{figure}[t]
  \centering
  \includegraphics[width=\linewidth]{static_pairing.png}
  \caption{Static clustering: long cross-cell edges reflect large gain disparity between paired users.}
  \label{fig:static}
\end{figure}

\subsection{Balanced Pairing}
To improve fairness, the sorted list is divided into two halves, pairing corresponding elements:
\[
\mathcal{P}_{\text{bal}}=\{(u_{(i)},u_{(i+N/2)})\mid i=1,\ldots,N/2\}.
\]
This configuration reduces excessive disparity between users while maintaining SIC feasibility.  
As shown in Fig.~\ref{fig:balanced}, balanced clustering yields shorter, more localized connections compared to static pairing, leading to improved user-rate uniformity and fairer power allocation.

\begin{figure}[t]
  \centering
  \includegraphics[width=\linewidth]{balanced_pairing.png}
  \caption{Balanced clustering: shorter intra-region edges moderate channel disparity and improve fairness.}
  \label{fig:balanced}
\end{figure}

\subsection{Blossom Matching (General Graph)}
A complete undirected graph is constructed where each vertex represents a user, and an edge $(i,j)$ exists if the pair satisfies the SIC condition.  
Each edge is weighted by the achievable sum rate $w_{ij}=R_{\mathrm{sum}}(i,j)$, computed from the SINR model.  
Edmonds' Blossom algorithm is then applied to determine the maximum-weight matching (MWM), ensuring global optimality with complexity $\mathcal{O}(N^3)$.  
As illustrated in Fig.~\ref{fig:blossom}, the Blossom clustering method adaptively forms pairs with optimal rate contribution while maintaining SIC feasibility across the network.  
Although it achieves the highest theoretical throughput, its computational cost grows rapidly with user density.

\begin{figure}[t]
  \centering
  \includegraphics[width=\linewidth]{blossom_pairing.png}
  \caption{Blossom clustering: global optimum pairing on a general graph with the highest computational complexity.}
  \label{fig:blossom}
\end{figure}

\subsection{Bipartite PF-Weighted Matching}
The user set is partitioned into weak ($U_W$) and strong ($U_S$) groups, either by a fixed median threshold or adaptively based on channel spread.  
Edges are allowed only across $(U_W,U_S)$ pairs that satisfy angular separation and SIC feasibility guards.  
Each feasible edge is assigned a proportional-fair (PF) weight:
\[
w_{ij} = \log(R_i+\epsilon) + \log(R_j+\epsilon),
\]
where $R_i$ and $R_j$ are achievable rates from the SINR model, and $\epsilon$ is a small offset to avoid $\log(0)$.  
The Hungarian or min-cost max-flow algorithm then produces an optimal bipartite matching in $\mathcal{O}(N^3)$ time.  
As shown in Fig.~\ref{fig:bipartite_pf}, the PF-weighted bipartite approach aligns naturally with the SIC structure of NOMA by pairing spatially separated weak–strong users, thereby achieving an excellent balance between fairness and total throughput.

\begin{figure}[t]
  \centering
  \includegraphics[width=\linewidth]{bipartite_pf_pairing.png}
  \caption{Bipartite PF clustering: aligns with NOMA SIC structure, offering high throughput and fairness with moderate complexity.}
  \label{fig:bipartite_pf}
\end{figure}

\subsection{OMA Fill and Throughput Accounting}
After matching, any unpaired users are served via orthogonal multiple access (OMA) on remaining sub-bands.  
Let $\mathcal{P}$ denote the set of NOMA pairs and $\mathcal{U}_{\text{OMA}}$ the set of unpaired users.  
Each user’s throughput is
\[
T = R \cdot \frac{B_{\text{unit}}}{10^6}, \qquad 
B_{\text{unit}} = \frac{B}{|\mathcal{P}| + |\mathcal{U}_{\text{OMA}}|},
\]
where $R$ is the user’s instantaneous rate (in bps/Hz) and $B_{\text{unit}}$ is the allocated bandwidth share.  
This ensures consistent accounting of aggregate system throughput across both NOMA and OMA users.
